% from ega4-2-fr.tex
\oldpage[IV-2]{124}

une $A_\mathfrak p$-algèbre finie; comme $S^{-1}\mathfrak p'$ est au-dessus
de l'idéal maximal $\mathfrak pA_\mathfrak p$ de $A_\mathfrak p$, c'est un
idéal maximal de $S^{-1}A'$; mais en vertu de l'hypothèse (ii) et de (5.6.3,
(i)), $A_\mathfrak p$ est universellement caténaire. On déduit donc de
(5.6.10) que
$\dim(A_\mathfrak p)=\dim(A'_{\mathfrak p'})=1$. Il résulte de là que l'on a
$A^{(1)}\subset A'$, et comme $A'$ est un $A$-module de type fini par
hypothèse, il en est de même de $A^{(1)}$ puisque $A$ est noethérien; la
proposition est donc démontrée dans le cas (ii).

\begin{remark}[5.11.3]
\label{IV.2.s5.11.3-fr}
Le fait que $A^{(1)}$ est une $A$-algèbre finie n'est plus nécessairement exact
si, dans l'hypothèse (ii) de (5.11.2), on suppose seulement que $A$ est
caténaire. Un exemple est donné par l'anneau local caténaire $A$ construit dans
(5.6.11), dont la clôture intégrale $A'$ (notée $B$ dans (5.6.11)) est une
$A$-algèbre finie; si $A^{(1)}$ était aussi une $A$-algèbre finie, comme il est
contenu dans le corps des fractions de $A$, il serait contenu dans $A'$. Mais
d'autre part, avec les notations de (5.6.11), tout idéal premier de hauteur 1
dans $A$ est de la forme $\mathfrak pA$, où $\mathfrak p$ est un idéal premier
de hauteur 1 dans $C$, et l'on a
$A_{\mathfrak pA}=C_\mathfrak p$; on sait (5.6.11) que
$C_\mathfrak p=E_{\mathfrak p'}$, où $\mathfrak p'$ est l'unique idéal
premier de $E$ au-dessus de $\mathfrak p$, donc $A_{\mathfrak pA}$ est
intégralement clos et contient par suite $A'$; par définition, on a donc
$A'\subset A^{(1)}$, et l'hypothèse que $A^{(1)}$ est une $A$-algèbre finie
entraînerait finalement $A^{(1)}=A'$. Mais cette conclusion est absurde, car un
des deux idéaux premiers de $A'$ au-dessus de l'idéal maximal $\mathfrak nA$ de
$A$ est de hauteur 1, alors que $\mathfrak nA$ est de hauteur 2, ce qui
contredirait (5.10.17, (iv)).
\end{remark}

\begin{corollary}[5.11.4]
\EGAReaderTarget{ega4.c04.cor.5-11-4}{IV.2.s5.11.4-fr}
Soient $X$ un préschéma localement noethérien, $Z$ une partie fermée de $X$,
$U=X-Z$, $i:U\to X$ l'injection canonique, $\mathcal F$ un
$\mathcal O_U$-Module cohérent. Pour que $i_*(\mathcal F)$ soit un
$\mathcal O_X$-Module cohérent, il est nécessaire que pour tout
$x\in\operatorname{Ass}(\mathcal F)$, on ait
\[
\operatorname{codim}(\overline{\{x\}}\cap Z,\overline{\{x\}})\geq2.
\]
Cette condition est suffisante dans chacun des deux cas suivants:
\begin{enumerate}
\item[(i)] Le préschéma $X$ est localement immersible dans un schéma régulier.
\item[(ii)] Le préschéma $X$ est universellement caténaire et universellement
japonais (ce qui signifie, par définition, que tout point $x\in X$ admet un
voisinage ouvert affine dont l'anneau est universellement japonais (0,
23.1.1)).
\end{enumerate}
\end{corollary}

On sait (I, 9.4.7) qu'il existe un sous-$\mathcal O_X$-Module cohérent
$\mathcal G$ de $i_*(\mathcal F)$ tel que $\mathcal G\vert U=\mathcal F$. On a
évidemment
$\operatorname{Ass}(\mathcal F)\subset\operatorname{Ass}(\mathcal G)
\subset\operatorname{Ass}(i_*(\mathcal F))$, et comme
\[
\operatorname{Ass}(i_*(\mathcal F))=\operatorname{Ass}(\mathcal F)
\]
(3.1.13), on a $\operatorname{Ass}(\mathcal G)=\operatorname{Ass}(\mathcal F)$;
il suffit alors d'appliquer (5.11.1) au $\mathcal O_X$-Module cohérent
$\mathcal G$, en notant que $i_*(\mathcal F)=\mathcal H^0_{X/Z}(\mathcal G)$
et que, lorsque $X$ est universellement caténaire et universellement japonais,
l'hypothèse (ii) de (5.11.1) est vérifiée par définition.

\begin{corollary}[5.11.5]
\label{IV.2.s5.11.5-fr}
Soient $X$ un préschéma localement noethérien, $Z$ une partie de $X$ stable par
spécialisation, $\mathcal F$ un $\mathcal O_X$-Module cohérent tel que l'on ait
$\operatorname{Ass}(\mathcal F)\subset X-Z$. Alors la condition:
\begin{enumerate}
\item[a)] $\mathcal H^0_{X/Z}(\mathcal F)$ est un $\mathcal O_X$-Module
cohérent
\end{enumerate}
implique la suivante:
\begin{enumerate}
\item[d)] Pour toute partie $T$ de $Z$ fermée dans $X$ (ou seulement pour une
famille filtrante

